\section{\gymname{}} 
\gymname{} is a hierarchical control framework built on the standard \textbf{Gymnasium} interface~\citep{brockman2016openai}. It binds a Low-Level Environment loop (a physics simulator or the real world) with a stateful Code Executor loop. 
Architecturally, \gymname{} preserves the native dynamics of underlying simulators, e.g., Robosuite~\citep{zhu2020robosuite}, LIBERO-PRO~\citep{liu2023libero, zhou2025liberopro}, and BEHAVIOR~\citep{li2024behavior1k}, while exposing them through a Read-Eval-Print Loop (REPL) paradigm tailored for coding agents. In CaP-Gym, a code environment ``turn" corresponds to one interaction between the coding agent and a specific robot task instance: the agent receives observations, generates a Python program, and the environment executes it to completion. The program may invoke multiple perception and control primitives, each of which can run the simulator or robot controller for multiple internal updates.

\subsection{Low-Level Perception and Control Primitives}
All computationally intensive perception and control primitives are implemented as stateless services~\cite{uvicorn}, enabling high-throughput parallel evaluation.

\textbf{Perception Primitives} 
Agents access perceptual data from the environment through modular perception primitives that abstract raw sensor data into structured semantic objects, e.g., SAM3~\cite{carion2025sam3} for language-conditioned segmentation and Molmo 2~\cite{clark2026molmo2} for open-vocabulary pointing, alongside standard vision libraries like OpenCV~\citep{opencv_library} and Open3D~\citep{zhou2018open3d}. 

\textbf{Control Primitives} Instead of directly emitting joint-space action commands, agents call motion planners or inverse kinematics solvers such as PyRoki~\cite{kim2025pyroki}. They can handle collision checking, reachability constraints, and action-space transformations, allowing agents to reason in a task-oriented Cartesian space while delegating execution feasibility to the controller. 
